% !TEX root = /home/eb/.vscode/extensions/qftrules.programme-colle/tmp/Exercice.tex


\begin{exo}[PC][1][devoir]{Décharge d'un conducteur dans l'air}
  Une boule de centre $\rm O$ et de rayon $R$, porte initialement la charge $Q_0$ uniformément répartie en volume, de densité $\rho$. Elle est abandonnée dans l'air, de conductivité électrique $\sigma$. On pose $Q(t)$ la charge électrique totale de la boule à l'instant $t$. On cherche à déterminer le champ électromagnétique $(\vv{E}(M,t),\vv{B}(M,t))$ en un point $\rm M$ de l'espace en dehors de la boule repéré en coordonnées sphériques $(r,\theta,\vphi)$.
  \begin{quest}
    \item  On suppose que la distribution de charges et de courants résultant de la décharge est isotrope à tout instant. Analyser les invariances et symétries de la distribution de charge $\rho$ et de courant~$\j$.
    \sol{
    Attention : ce sont $\rho$ et $\j$ qui définissent les plans de symétrie ou d’antisymétrie. Ça n’a donc pas de sens de trouver un plan de symétrie, puis de dire que $\j$ est orthogonal ou compris dans ce plan.\\
  
  
    Symétrie sphérique : $\rho(M) = \rho(r)$. Par propriété d’isotropie, on a ensuite
  \eqbox{\j(M) = j(r)\er.}
    }
  \item  Rappeler la loi d'Ohm locale. En déduire la forme du champ électrique $\vv{E}(M,t)$ à l'extérieur de la boule, puis justifier que $\vv{B}(M,t)=0$.
  \sol{\begin{liste}
    \item Champ électrique. Loi d'Ohm : $\E(M) = \sigma \j(M) = \sigma j(r) \er$, soit
  \eqbox{\E(M) = E(r) \er.}
    \item Champ magnétique. Les plans $(M,\et,\er)$ et $(M,\ep,\er)$ sont des plans d'antisymétrie de la distribution de courant, donc $\B(M) = B(r) \er$. Mais $\rot \B = -\mu_0 \epz \dr{\E}{t}$ est aussi porté par $\er$. La seule solution est que
  \eqbox{\B=\Vec{0}.}
  \end{liste}
  
  }
  \item En utilisant le théorème de Gauss et la conservation de la charge, établir l'équation différentielle vérifiée par la charge $Q(t)$. La résoudre et expliquer pourquoi la boule se déchargera d'autant plus rapidement que l'air est humide.
  \sol{
  Deux obtions : Maxwell-Ampère avec $\rot \B = 0$, ou bien théorème de Gauss avec conservation de la charge. On montre ici la deuxième méthode.
  
  
  Conservation de la charge : $\dd{Q}{t} =- \oiint_S \j \ps \dS = -4\pi R^2 j(R)$. Or d'après le théorème de Gauss,
  $$
  4\pi R^2 E(R) = 4\pi R^2 j(R)/\sigma  = \frac{Q(t)}{\epz}.
  $$
  donc
  $$
  \dd{Q}{t} = - \frac{\sigma}{ \epz} Q
  ,
  $$
  Soit
  \eqbox{\dd{Q}{t}+\frac{Q}{\tau}=0}
  avec $\tau = \frac{\epz}{\sigma}$. Dans un air humide, $\sigma$ est plus élevé car l’air humide est meilleur conducteur, mais $\epz$ est un peu plus élevé aussi. En pratique, c’est $\sigma$ qui gagne et $\tau$ est donc plus petit.
  }
  \item Vérifier que les expressions satisfont l’équation locale de Poynting.
  \sol{
  Ici $\B=\Vec{0}$ et donc le vecteur de Poynting $\Poy = \frac{\E \pv \B}{\mu_0}$ est nul,
  $$
  \Poy = \Vec{0}.
  $$
  L’équation de conservatior de l’énergie électromagnétique s’écrit alors
  \eqbox{\dr{u_{\rm em}}{t} = -\j \ps \E.}
  Le premier terme vaut en $r=R$,
  $$
  \dr{u_{\rm em}}{t} = \dr{}{t}\left( \frac{1}{2}\epz E^2\right)=  \dr{}{t}\left( \frac{Q^2}{32\epz\pi^2 R^4}\right) =  \frac{Q}{16\epz\pi^2 R^4} \dot{Q} = -\frac{Qj}{4\epz\pi R^2},
  $$
  de par la conservation de la charge.
  Le deuxième terme vaut en $r=R$,
  $$
  -\j \ps \E = -\frac{Qj}{4\epz\pi R^2} = \dr{u_{\rm em}}{t} .
  $$
  }
  \end{quest}
  \end{exo}
  %%%%%%%%%%%%%%%%%%%%%%%%%%%%%%%%%%%%%%%%%%%%%%%%%%%%%%%

  
%%%%%%%%%%%%%%%%%%%%%%%%%%%%%%%%%%%%%%%%%%%%%%%%%%%%%%%%%%
\begin{exo}[PC][1][colle]{Chauffage par induction}
Un cylindre métallique de rayon $a$, de hauteur $h$ et de conductivité $\sigma=\SI{5e7}{S.m^{-1}}$ est placé à l'intérieur d'une bobine supposée infinie suivant son axe $(Oz)$. Cette bobine est parcourue par un courant variable $\frac{\omega}{2\pi}=f=\SI{50}{Hz}$ et génère un champ magnétique uniforme mais dépendant du temps $\vv{B}=B_{0}\cos(\omega t)\,\vv{e_z}$.
\begin{quest}
\item Déterminer l'expression du champ $\vv{E}$ en tout point du conducteur, en négligeant devant $\vv{B}$ le champ magnétique induit généré par le courant induit. En déduire l'expression de la densité volumique des courants induits $\vv{j}$, dits courants de Foucault.
\item Déterminer la moyenne sur une période $T=\frac{2\pi}{\omega}$ de la densité volumique de puissance dissipée à une distance $r$ de l'axe par effet Joule, puis la puissance moyenne totale dissipée sur tout le conducteur. Dans quel but cette perte est-elle ici provoquée ?
\item À l'aide du théorème d'Ampère, déterminer l'expression du champ magnétique induit $\vv{B_{\rm{ind}}}$ créé par les courants induits de la question 1. On admettra que $\vv{B}_{\rm{ind}}(r=a,t)=0$ à la surface du conducteur. L'hypothèse $||\vv{B}||\gg||\vv{B_{\rm{ind}}}||$ est-elle vérifiée si $a\sim\SI{10}{cm}$ ?
\item Sans faire l'approximation $||\vv{B}||\gg||\vv{B_{\rm{ind}}}||$, déterminer à quelle équation satisfait le champ magnétique total $\vv{B_{\rm{tot}}}$ dans l'ARQS magnétique. À quel type de solution faut-il s'attendre ? Commenter.
\end{quest}
\end{exo}
%%%%%%%%%%%%%%%%%%%%%%%%%%%%%%%%%%%%%%%%%%%%%%%%%%%%%%%%%%


%%%%%%%%%%%%%%%%%%%%%%%%%%%%%%%%%%%
\begin{exo}[PC][1][colle]{Freinage par induction}
  Une spire carrée de coté $a$, de masse $m$, tombe dans le champ de pesanteur $\Vec{g}=g \ex$. Dans le demi-espace $x>0$, règne le champ magnétique uniforme et permanent $\B=B_0 \ez$. À l'instant $t=0$, la spire se trouve dans la situation représentée sur la figure ci-dessous, sa vitesse est $\Vec{v}=v_0 \ex$, son côté inférieur est en $x=0$.
\begin{center}
  \includegraphics[width=0.5\columnwidth]{carre_induction.png}
\end{center}
\begin{quest}
  \item
   Montrer que le mouvement ultérieur de la spire reste une translation verticale selon l'axe $O x$.
   \item Soit $R$ la résistance de la spire. Déterminer la vitesse $v(t)$ de la spire.
   \item La spire a maintenant une résistance nulle (spire supraconductrice) et on note $L$ son inductance propre. Reprendre l'étude précédente et préciser la condition d'oscillation de la spire.
\end{quest}
\end{exo}
%%%%%%%%%%%%%%%%%%%%%%%%%%%%%%%%%%%


%%%%%%%%%%%%%%%%%%%%%%%%%%%%%%%%%%%%%%%%%%%%
\begin{exo}[PC][3][TD]{Énergétique d’un fil électrique}
  \Notions{\item Géométrie cylindrique \task Vecteur de Poynting \task Électrostatique \task Magnétostatique}
  On étudie un fil électrique soumis à une différence de potentiel continue dans lequel circule un courant constant d’intensité $I$. Le fil est constitué d’un conducteur ohmique de conductivité électrique $\gamma$. On néglige les effets de bord, de sorte à assimiler le fil à un cylindre infini de rayon $a$ et d’axe $(Oz)$. On note $\j = j\,\ez$ le vecteur densité de courant dans le fil, que l’on suppose uniforme et stationnaire.

\begin{quest}
\item  Déterminer l'expression du champ électrique $\E(M)$ à l’intérieur du fil en fonction de $I$.
\sol{
Loi d’Ohm locale :
$$
\E(r<a) = \frac{\j}{\gamma} = \frac{j}{\gamma} \ez.
$$
}
\item Déterminer l’expression du champ magnétique $\B(M)$ en à l'intérieur du fil en fonction de $I$.
\item En déduire le vecteur de  Poynting $\Poy(M)$ à l’intérieur du fil.  Représenter les vecteurs $\E$, $\B$ et $\Poy$ sur un schéma.  Interpréter en réalisant un bilan d’énergie électromagnétique. L'énergie électromagnétique est-elle transportée par les porteurs de charge à l'intérieur du fil ?
\sol{
$\B = \frac{\mu_0 j r}{2} \et$ Donc
$$
\Poy= \frac{1}{\mu_0} \E \pv \B = - \frac{j^2 r}{2 \gamma} \er.
$$
Le flux du vecteur de Poyting en $r=a$ vers l’intérieur du fil vaut
$$
\iint_S \Poy \ps \dS = \frac{j^2 a}{2 \gamma} \times (2\pi a h) = \frac{j^2}{\gamma} (\pi a^2 h).
$$
Ce flux compense bien la puissance perdue par effet Joule dans le cylindre
$$
\iiint \j \ps \E \d V = \frac{j^2}{\gamma} V = \frac{j^2}{\gamma} (\pi a^2 h).
$$
L’énergie électromagnétique $U_{\rm em}$ contenue dans ce cylindre est constante, comme il se doit en stationnaire.
}
\end{quest}
 On admet qu'un conducteur électrique à l’équilibre est chargé uniformément en surface. Pour un conducteur parfait, la densité surfacique de charge  $\sigma$ à la surface du fil est uniforme, de même que le potentiel électrostatique $U$. On suppose que le fil constitue l'âme d'un câble coaxial entouré d'une gaine de rayon $R \gg a$ et portée à un potentiel nul.
\begin{quest}
\item Analyser les invariances et symétries de cette distribution, puis déterminer le champ électrique $\E(M)$ et le champ magnétique $\B(M)$ à l’extérieur du fil. Déterminer la charge surfacique $\sigma$ en fonction du potentiel~$U$.
\sol{
$\E(r>a,z) = \frac{\sigma a}{r \epz} \er$.
La tension $U$ entre la surface du fil et la distance $R$ est donnée par
$$
V(r=a) - V(r=R) = U = \int_R^a \E \ps \er \d r = \frac{\sigma a}{\epz} \int_R^a \frac{\d r}{r} = \frac{\sigma a}{\epz} \ln\left( \frac{a}{R}\right),
$$
et donc 
\eqbox{
\sigma = \frac{\epz U}{a \ln\left( \frac{a}{R}\right)}.
}
}
\item En déduire le vecteur de Poynting~$\Poy(M)$ à l'extérieur du fil. Représenter les vecteurs $\E$, $\B$ et $\Poy$ sur un schéma et interpréter. Calculer la puissance circulant autour du fil conducteur. 
%Déterminer la puissance totale transférée par le champ électromagnétique le long du fil.
 Commenter.
%pourquoi ce ne sont pas les porteurs de charge mais le champ électromagnétique qui apporte directement l’énerie de la source de tension au récepteur.
\sol{
On sait que
$\B(r>a) = \frac{\mu_0 I}{2\pi r} \et$.
Donc
$$
\Poy = \frac{1}{\mu_0} \E \pv \B  = \frac{\sigma I }{2\pi r^2 \epz} \ez.
$$
La puissance totale vaut
$$
\mc{P} = \int \Poy \ps \d S \ez = \int_a^R \frac{\sigma I }{2\pi r^2 \epz}  \times 2\pi r \d r = \frac{\sigma I }{\epz}  \int_{a}^R \frac{\d r}{r} = \frac{\sigma I }{ \epz} \ln\left( \frac{R}{a}\right) = {I U}.
$$
}
\end{quest}
Dans un conducteur réel, le potentiel à l'intérieur du fil conducteur $V(r<a,z)$ est une fonction affine de la position $z$ le long du fil. Il en va donc de même de la charge surfacique $\sigma(z)$. De plus, le champ extérieur possède une composante axiale $E_z(r>a)$. On cherche l'expression du potentiel $V(r,z)$ à l'extérieur du fil sous la forme 
\begin{equation}
V(r,z) = A \ln(r) + B z + C z \ln(r) + D,
\end{equation}
sachant que le potentiel vaut $V(r=a,z=0) = 0$ en amont du fil et $V(r=a,z=L) = U$ en aval du fil de longueur $L$, avec toujours $V(r=R,z) = 0$ à la surface de la gaine.
\begin{quest}
  \item	Vérifier que l'expression proposée satisfait à l'équation de Laplace. En utilisant les conditions aux limites, déterminer les coefficients $A$, $B$, $C$ et $D$ en fonction de $U$, $a$ et $R$.
  \item En déduire l'expression du champ électrique $\E(r,z)$ à l'extérieur du fil et de la charge surfacique $\sigma(z)$.
  \item On représente ci-dessous une résolution numérique des équations de Maxwell pour un circuit constitué d'un générateur de tension $U$ alimentant deux résistances $R$ à travers un fil conducteur. Commenter.
\end{quest}
\begin{fig}[label]
  \centering
  \begin{subfig}[label][0.4]
    \includegraphics[height=\columnwidth]{electric_circuit_E.png}
  \end{subfig}\hspace{0.5cm}
  \begin{subfig}[label][0.4]
    \includegraphics[height=\columnwidth]{electric_circuit_pi.png} 
  \end{subfig}
  \capt{(a) Champ électrique et équipotentielle. (b) Vecteur de Poynting. Source : S. Macjen, R.K. Haaland and S.C. Dudley, Am. J. Phys. 68,  857 (2000).}
\end{fig}
\end{exo}
%%%%%%%%%%%%%%%%%%%%%%%%%%%%%%%%%%%%%%%%%%%%


%%%%%%%%%%%%%%%%%%%%%%%%%%%%%%%%%%%%%%%%%%%%%%%%%%%%%%%
\begin{exo}[PC][1][TD]{ARQS magnétique dans un conducteur ohmique}
  \Notions{\item ARQS \task Conducteur ohmique}
  Un conducteur ohmique de conductivité électrique $\sigma$ est soumis à un champ électromagnétique $(\E,\B)$. On note $\rho(M,t)$ la densité volumique totale de charge, comprenant électrons et cations, et $\j(M,t)$ la densité de courant totale dans le conducteur.
\begin{quest}
  \item Montrer que l’équation de conservation de la charge impose un décroissance de $\rho$ sur un temps typique $\tau_n$ que l’on déterminera. Comparer ce temps typique de retour à la neutralité $\tau_n$ au temps typique des collisions des électrons avec la réseau cristallin $\tau$ pour le cuivre, de conductivité $\sigma = \SI{6.0e7}{S.m^{-1}}$.
  \sol{
  L’équation de conservation de la charge s’écrit
  $$
  \dd{\rho}{t} + \dive \j = 0.
  $$
  L'équation de Maxwell-Gauss s'écrit
  $$
  \dive \E = \frac{\rho}{\epz},
  $$
  et la loi d'Ohm, valable dans l'AQRS, s'écrit
  $$
  \j = \sigma \E.
  $$
  On en déduit que
  $$
  \dd{\rho}{t} + \dive \left(\sigma \E\right) = 0 \Rar \dd{\rho}{t} + \frac{\sigma}{\epz} \rho = 0.
  $$
  On peut introduire le temps typique $\tau_n = \frac{\epz}{\sigma} = \frac{m\epz }{ne^2 \tau} \simeq \SI{e-18}{s}$. Ce temps est bien plus petit que le temps typique des collisions des électrons avec la réseau cristallin $\tau \simeq \SI{e-14}{s}$. Dans l'ARQS, on peut donc supposer que le métal est neutre électriquement.
  
  }
  \item Rappeler les conditions de l’ARQS magnétique, et montrer que cette approximation est valable aux fréquences $f$ faibles, telles que $f\tau_n\ll 1$. En déduire que l'ARQS dans un conducteur ohmique assure aussi l'ARQS magnétique.
\end{quest}
\end{exo}
%%%%%%%%%%%%%%%%%%%%%%%%%%%%%%%%%%%%%%%%%%%%%%%%%%%%%%%

%%%%%%%%%%%%%%%%%%%%%%%%%%%%%%%%%%%%%%%%%%%%%%%%%%%%%%%%%%
\begin{exo}[PC][2][TD]{Effet Hall et magnétorésistance}
  \Notions{\item Effet Hall \task Résistance électrique \task Loi d'Ohm}
On considère un conducteur ohmique de conductivité $\sigma$ soumis à un champ magnétique transverse $\vv{B}=B\, \ez$ et un champ électrique axial $\E=E\, \ex$ tous deux uniformes et stationnaires. On souhaite déterminer la loi d'Ohm locale reliant $\j$ en $\E$ en présence de ce champ magnétique transverse~$\B$.
\begin{quest}
\item Appliquer le PFD à un électron de conduction, de charge $-e$, masse $m$, et vitesse $\Vec{v}$, et projeter dans le plan $(x,y)$.  Résoudre le système d'équations différentielle obtenu en régime stationnaire.
\sol{
  \udl{Système} : $\{$électron$\}$.


  \udl{BDF} : \begin{itemize}
        \item force électrique : $\Vec{F}_E = -e \Vec{E}$ ;
        \item force magnétique : $\Vec{F}_B = -e \Vec{v} \pv \B$ ;
        \item force de frottement phénoménologique : $\Vec{F} = -\frac{m}{\tau} \Vec{v}$.
  \end{itemize}
  \udl{PFD} :
$$
m \dd{\Vec{v}}{t} = -e \Vec{E}- e \Vec{v} \pv \B - \frac{m}{\tau} \Vec{v}
$$
En projection, on trouve
$$
\begin{cases}
  m \displaystyle\dd{v_x}{t} & = -eE - \dfrac{m}{\tau} v_x -ev_y B, \\[3mm]
  m \displaystyle\dd{v_y}{t} & = 0 - \dfrac{m}{\tau} v_y +ev_x B. \\
\end{cases}
$$
En régime stationnaire,  la vitesse est constante et l’on trouve
$$
\begin{cases}
  Ee & = - \dfrac{m}{\tau} v_x -ev_y B, \\[3mm]
  0 & = - \dfrac{m}{\tau} v_y +ev_x B. \\
\end{cases}
$$
La seconde relation donne 
$$
v_y = \frac{eB\tau}{m} v_x.
$$
La première relation devient
$$ Ee  = - \frac{m}{\tau} v_x  - e^2 \frac{\tau}{m} B^2 v_x = - \frac{m}{\tau} v_x \left[ 1+\left( \frac{eB\tau}{m} \right)^2 \right]$$ et ainsi
\eqbox{
  \begin{cases}
    v_x & = -\dfrac{e\tau}{m} \dfrac{E}{1+\left(\dfrac{e \tau B}{m }\right)^2}, \\[10mm]
    v_y & = -\left(\dfrac{e\tau}{m}\right)^2 \dfrac{E B}{1+\left(\dfrac{e \tau B}{m }\right)^2}.\\
  \end{cases}
  }
}

\item En déduire la composante $j_x$ de la densité de courant en fonction de $E$, $\sigma_0$, $B$ et la constante de Hall du matériau $R_H$. Montrer que la résistance électrique longitudinale $R$ d'un tube conducteur de longueur $L$ et de section $S$, a pour expression 
\begin{equation}
  R = R_0 \left[ 1+(\sigma R_H B)^2 \right],
\end{equation}
où $R_0$ désigne la résistance du tube en l'absence de champ magnétique, à exprimer en fonction de $\sigma_0$, $L$ et $S$.
Justifier l'appellation de magnétorésistance.  De quel facteur est augmenté la résistance pour du cuivre et du germanium, soumis à un champ magnétique $B=\SI{1}{T}$ ? \AN~?
\sol{
  Le vecteur densité de courant de charges s’écrit
$$
\j = n q \Vec{v},
$$
où $q$ est la charge des porteurs mobiles et $\Vec{v}$ leur vitesse moyenne.
On a donc
$$
j_x = -nev_x = \frac{ne^2\tau}{m} \frac{E}{1+\left(\frac{e \tau B}{m }\right)^2} = \frac{\sigma_0}{1+(\sigma_0 R_H B)^2} E  = \sigma_{\parallel} E,
$$
avec $R_H=1/ne$ la constante de Hall, et $\sigma_0 = ne^2 \tau/m$ la conductivité électrique longitudinale en champ nul.
La résistance longitudinale en présence de champ magnétique vaut donc
$$
R = \frac{L}{\sigma_{\parallel} S} = \frac{L}{\sigma_0  S} \left[ 1+(\sigma_0 R_H B)^2 \right],
$$
soit 
\eqbox{
R = R_0 \left[ 1+(\sigma_0 R_H B)^2 \right],
}
où $R_0 = L/(\sigma_0 S)$ est la résistance du matériau en l'absence de champ magnétique.
La résistance augmente avec le champ magnétique $B$. On parle de magnétorésistance. Pour du cuivre, de conductivité $\sigma_0 = \SI{6.0e7}{S.m^{-1}}$, soumis à un champ magnétique $B=\SI{1}{T}$, on trouve
\eqbox{
\frac{R}{R_0} = 1 + (\sigma_0 R_H B)^2
}
}
\item Déterminer la composante $j_y$ de la densité de courant en fonction de $E$, $\sigma$, $B$ et la constante de Hall du matériau $R_H$. Interpréter l'existence de ce courant transverse en terme d'effet Hall. Déterminer la conductivité transverse  $\sigma_{H}$, aussi appelée conductivité de Hall. \AN~?
\sol{
  Le courant transverse vaut
$$
j_y = -nev_y = ne \left( \dfrac{e\tau}{m}\right)^2 \dfrac{E B}{1+\left(\frac{e \tau B}{m }\right)^2}  = \dfrac{\sigma_0^2 R_H B}{1+(\sigma_0 R_H B)^2} E.
$$
On peut généraliser la loi d'Ohm en présence d'un champ magnétique en reliant $\j$ à $\E$ par une matrice conductivité, notée $[\sigma]$. Cette matrice doit être symétrique. À partir des coefficients trouvés précédemment, on peut écrire
% $$
\eqbox{
  \j =  [\sigma] \E = \frac{\sigma_0}{1+(\sigma_0 R_H B)^2} 
  \begin{pmatrix}
    1 &  \sigma_0 R_H B\\
    \sigma_0 R_H B   & 1 \\
  \end{pmatrix} \E
  }
% $$
}
\end{quest}


\begin{tabular}{lll}
  \textsc{Données} :  & conductivité électrique du cuivre : & $\sigma_0 = \SI{6.0e7}{S.m^{-1}}$ \\
  &constante de Hall du cuivre : & $R_H = \SI{1.0e-10}{m^3.C^{-1}}$ \\
  &conductivité électrique du germanium : & $\sigma_0 = \SI{2.0e-3}{S.m^{-1}}$ \\
  &constante de Hall du germanium : & $R_H = \SI{1.0e-6}{m^3.C^{-1}}$ \\
  &charge élémentaire : & $e = \SI{1.6e-19}{C}$ \\
  &masse de l'électron : & $m = \SI{9.1e-31}{kg}$ \\
\end{tabular}
\end{exo}
%%%%%%%%%%%%%%%%%%%%%%%%%%%%%%%%%%%%%%%%%%%%%%%%%%%%%%%%%%

%%%%%%%%%%%%%%%%%%%%%%%%%%%%%%%%%%%%%%%%%%%%%%%%%%%%%%%%%%
\begin{exo}[PC][1][colle]{Magnétorésistance}
En présence d'un champ $\vv{B}=B\,\vv{e_z}$ uniforme et stationnaire, on admet que la loi d'Ohm dans un métal est remplacée par la relation $\vv{j}+\alpha\vv{j}\pv\vv{B}=\sigma\vv{E}$ où $\alpha$ est une constante positive.
\begin{quest}
\item Interpréter la relation précédente en faisant le lien avec l'effet Hall. En déduire l'expression de $\alpha$ en fonction de la conductivité électrique $\sigma$, la densité volumique de porteurs de charges $n$ et la charge élémentaire $e$.
\sol{
La force magnétique s'écrit
}
\item Un métal remplit sur une hauteur $h$ l'espace entre deux cylindres coaxiaux de rayons $r_1$ et $r_2>r_1$. Ces cylindres sont des équipotentielles portées à $V_1=0$ et $V_2=U$. Le tout est plongé dans un champ magnétique $\vv{B}=B\,\vv{e_z}$ uniforme et stationnaire.
\begin{enumerate}
\item En analysant les symétries et invariances, montrer que l'expression de $\vv{E}$ est particulièrement simple. Déterminer alors les expressions des composantes de $\vv{j}$ en fonction de celle(s) de $\vv{E}$ ainsi que de $\sigma,\alpha$ et $B$.
\item Calculer la résistance électrique du conducteur en régime stationnaire. Commenter.
\end{enumerate}
\end{quest}
\end{exo}
%%%%%%%%%%%%%%%%%%%%%%%%%%%%%%%%%%%%%%%%%%%%%%%%%%%%%%%%%%

%%%%%%%%%%%%%%%%%%%%%%%%%%%%%%%%%%%%%%%%%%%%%%%%%%%%%%%%%%
\begin{exo}[PC][1][colle]{Transport de courant dans un fil}
Considérons une portion de fil de cuivre soumise à une tension harmonique de pulsation $\omega$ à ses bornes. On modélise le fil par un cylindre infini de rayon $a$ et d'axe $(Oz)$. On se place dans le cadre de l'ARQS de sorte que la conductivité $\sigma_0 = \SI{6e7}{S.m^{-1}}$ est réelle. On cherche la densité volumique de courant imposée par le générateur sous la forme $\udl{\j}(M,t)= \udl{j}(r)\,\e^{i\omega t}\,\ez$ en coordonnées cylindriques ($r,\theta,z$).
\begin{quest}
\item Justifier l'expression proposée pour $\udl{\j}$.
\item Montrer que dans un régime basse fréquence a définir, le champ électrique $\udl{\E}$ satisfait à l'équation
$$
\Delta \udl{\E} = \mu_0\sigma_0\dr{\udl{\E}}{t}.
$$
Justifier que la densité de courant $\udl{\j}$ satisfait à la même équation. De quelle type d’équation s’agit-il ?
\item On écrit $\udl{j}(r) = j_0 \exp(-\frac{1+i}{\delta}(a-r))$,
où $j_0$ désigne l'amplitude de la densité volumique de courant en périphérie du cylindre.
Quelle est le nom usuel et la dimension physique de $\delta$ ?
On rappelle l'expression du Laplacien en coordonnées cylindriques pour une fonction qui ne dépend que de $r$, $\Delta A=\frac{1}{r}\frac{\d}{\d r}\left(r\frac{\d A}{\d r}\right)$. Exprimer $\delta$ dans le cas où $r\gg\delta$ en fonction de $\mu_0,\sigma_0$ et $\omega$. Interpréter physiquement la forme de $\udl{j}$.
\item Évaluer numériquement $\delta$ pour une fréquence de \SI{50}{Hz} et de \SI{100}{MHz}. Pour un fil de section \SI{1}{mm^2} l'effet précédent est-il observable dans les deux cas ? Comparer la résistance du fil à ces deux fréquences et proposer une solution technologique permettant l'utilisation à haute fréquence.
\end{quest}
\end{exo}
%%%%%%%%%%%%%%%%%%%%%%%%%%%%%%%%%%%%%%%%%%%%%%%%%%%%%%%%%%
